\section{Direkt Demokrácia}

Az előző fejezetek élesen kritizálták a jelenlegi rendszert, de nem ajánlottak alternatívákat.
Azonban a demokráciának rengeteg formája létezik. Az, hogy ezek nem ismertek, nem meglepő egy
olyan oktatási rendszerben, amely igyekszik megtanítani a fiatalokat, mit gondoljanak, ahelyett
hogy azt tanítaná, hogyan gondolkodjanak.

A teljes átláthatóság érdekében ebben a fejezetben felsoroljuk ezeket az alternatívákat, majd
röviden megvitatjuk, hogy miért nem ezek mellett érvel ez a cikk. Ezután bemutatásra kerül a
direkt demokrácia egy formája, végezetül pedig útmutatást biztosítunk ezen rendszer megvalósításához
a gyakorlatban. Az olvasóra bízzuk, hogy elfogadja az érveinket, vagy a felsorolt alternatívák közül
egy másikat választ, és dolgozik azon rendszer implementálásán.

A következőkben felsorolunk néhány népszerű demokratikus alternatívát. Habár ezek a rendszerek megoldást
kínálnak megannyi problémára a képviseleti demokráciában, véleményünk szerint egyik sem ideális választás
a jelenlegi helyzetben. Néhány esetben ennek oka, hogy a rendszer túlságosan különbözik a jelenlegitől,
emiatt a megvalósítása nehézkes lenne. Más esetben az alternatíva azt a veszélyes hitet népszerűsíti,
hogy elég valaki másra bízni életünk legfontosabb kérdéseit; ha jól választunk, akkor jó döntések fognak
születni. Véleményünk szerint ez a legveszélyesebb gondolat a modern világban. Az ötlet, hogy egy idegen
--- legyen akármilyen képzett --- képes mindenki számára megfelelő döntést hozni, egy abszurdum, de
minimum oka a világ jelenlegi helyzetének. A jó képviselet egyetlen garanciája az önképviselet:

\begin{itemize}
	\item \textbf{Lotokrácia / Sorshúzásos Demokrácia} \cite{lottery_dem_1, lottery_dem_2, lottery_dem_3}:
	      választások helyett sorsolással választja ki a képviselőket. Tehát a törvényhozást egy véletlenszerűen
	      kiválasztott polgári csoport végzi. Ezt a csoportot néhány évente újrasorsolják. Előnye, hogy sokkal
	      jobban reprezentálja a társadalmat, hiszen nem csak azok kerülhetnek döntési pozícióba, akiknek ideje,
	      pénze és erőforrása engedi. Továbbá csökkenti a korrupciót, és megakadályozza, hogy a politika az
	      elitek játékává váljon.

	\item \textbf{Folyékony Demokrácia} \cite{liquid_dem_1, liquid_dem_2}: esetében a polgárok szavazhatnak
	      közvetlenül előterjesztésekre, vagy átruházhatják szavazati jogukat egy általuk megbízhatónak vélt félre.
	      Ez egy olyan rugalmas rendszert hoz létre, ahol a befolyás folyamatosan áramlik a legmegbízhatóbb emberek
	      felé. Ez a demokrácia igyekszik lépést tartani az állandóan változó világgal, és lehetővé teszi, hogy a
	      polgárok dinamikusan fektessék különböző szakértőkbe a bizalmukat, attól függően, hogy kiben bíznak az
	      adott előterjesztéssel kapcsolatban a legjobban.

	\item \textbf{Modern Közvetlen Demokrácia} \cite{direct_dem_1}: a polgárok a törvényhozást közvetlenül
	      végzik modern eszközök, mint az internet segítségével. Nincsenek választások vagy közvetítők; a demokrácia
	      folytonos, és minden döntés a többség akaratával születik.

\end{itemize}

A fentebb említett definíciókból látszik, hogy a közvetlen demokrácia nem egyetlen rendszerre utal,
hanem sokkal inkább egy gyűjtőfogalom, amely különböző demokratikus rendszereket foglal magába. A
következő fejezetben a közvetlen demokrácia egy formája mellett fogunk érvelni, amely egyszerre
megvalósítható és orvosolni tudja a jelenlegi rendszer problémáit.


\subsection{Fakultatív Modern Közvetlen Demokrácia}

A fakultatív modern közvetlen demokrácia egy olyan rendszer, amelyben a polgároknak lehetőségük van
minden előterjesztésre szavazni és sajátot alkotni, azonban konszenzus hiányában a döntéshozás
a megválasztott képviselők feladata. A rendszer "fakultatív", hiszen nem hárítja a polgárokra
a törvényhozás terhét, csupán teret ad a részvételhez. Azonban "közvetlen" abban az értelemben,
hogy a többségi akarat nem másítható meg a képvisleők által, legyenek ezek bármilyen pozícióban.

A gyakorlatban egy telefonos applikáció használható arra, hogy a polgárok gyorsan kifejezzék
véleményüket egy adott előterjesztéssel kapcsolatban. Fontos, hogy ez az applikáció a kormány
által támogatott legyen, de egy non-profit szervezet végezze a fejlesztését, és hogy a programkód
teljesen nyílt forrású legyen, ezzel garantálva a maximális átláthatóságot.

Ennek a rendszernek megannyi előnye lenne, itt van néhány a teljesség igénye nélkül:
\begin{itemize}
	\item Növeli az állampolgárok közvetlen részvételét a döntéshozatalban.

	\item Erősíti a demokratikus legitimitást, mivel a döntések közvetlenül a nép akaratát tükrözik.

	\item Csökkenti a politikai elit és a közvetítők szerepét.

	\item Fokozza az állampolgári felelősségérzetet és politikai tudatosságot.

	\item Lehetővé teszi a gyorsabb reakciót a társadalmi problémákra.

	\item Növeli az átláthatóságot a döntéshozatali folyamatban.

	\item Elősegíti a közérdek érvényesülését a pártérdekekkel szemben.

	\item Minimalizálja a korrupciót, hiszen a többségi akarat gátat szabhat a nem legitim intézkedéseknek.

	\item Segíti a képviselőket, mert valós idejű és pontos információt biztosít a népakaratról.

\end{itemize}

Azonban megannyi kritika is felmerülhet a rendszerrel kapcsolatban, ezért fontos hogy a leggyakoribbakat
megválaszoljuk:

\begin{enumerate}

	\item \textbf{Az embereknek nincs idejük a politikára.} Szóval vezessük be a 32-órás munkahetet változatlan
	      bérekkel. Ezáltal ez embereknek több idejük lesz az életüket szabályozó döntésekre. A 32-órás munkahét egy
	      tökéletes példa valamire, ami sosem kerülhet bevezetésre a jelenlegi rendszerben. Annak ellenére, hogy
	      rengeteg tanulmány kimutatta\cite{32_hours_ww_1, 32_hours_ww_2}, mennyire előnyös a munkavállalók számára és
	      a cégek ahol kiprobálták megvoltak elégedve az eredményekkel. Azonban a képviselők vagy maguk is cégtulajdonosok
	      vagy közel állnak több cégvezetőhöz a kampányuk vagy a családjuk révén. Angliában az újonnan választott
	      munkáspárt megpróbálta bevezetni a 4 napos munkahetet, de megtartották a 40 órát, így az embereknek 10 órát
	      kell naponta dolgozniuk. Úgy látszik a hegytetőről mindig úgy tűnik, hogy a hangyák nem dolgoznak elég
	      szorgalmasan.

	\item \textbf{Az átlag ember túl ostoba, hogy megértse a politikát és helyes döntést hozzon.} Sajnos az
	      átlagember a legjobb amink van. A képviselők is csak átlag emberek. Gyakran rosszabbak is az átlag embernél.
	      A legtöbbjük multimilliomos és már megvan mindenük, amire csak szükségük lehet. Sokuk felsőbb társadalmi
		  körökbe nőtt fel és kapcsolatok révén jutott lehetőségekhez. Megannyiuknak soha nem kellett a piacon
		  megméretetniük magukat, ellenben azon munkavállalókkal akiket elméletben képviselnek. 

	\item \textbf{Néhány őrült rossz törvényeket fog hozni.} Ez miben lenne más, mint a mostani rendszer?
	      A közvetlen demokráciában is hozhatnak rossz döntéseket, de ezeket a többségnek mindenképpen el kell
	      fogadnia. Azonban a jelenlegi rendszerrel ellentétben, nem kell majd négy évet várnunk, hogy meggondoljuk
	      magunkat. Egy láthatóan nem működő törvény gyorsan módosítható vagy visszavonható lesz ebben az új rendszerben.
	      Továbbá a döntés hozói és elszenvedői ugyan azok lesznek, ami sokat segíthet egy rossz döntés visszavonásában.

	\item \textbf{Soha nem fogunk majd semmit megvalósítani.} Mit tudunk megvalósítani most? Leszámítva a
	      közszolgáltatások költségvetésének csökkentését és a politikusok fizetésének megemelését? A második
	      világháború után a kormány egész lakóparkokat épített, és a közszolgáltatások működtek. Mi volt a
	      legnagyobb projekt amit a kormány az utóbbi 25 évben véghez vitt? A képvisleők mindig találnak időt,
	      hogy magukon segítsenek, de valahogy az egészségügy rendbehozatalára mindig kell még négy év. Az új
	      rendszerben sokkal gyorsabban fogunk haladni, mert könnyebb meggyőzni 1 millió embert valamiről ami
	      előnyös számukra, mint 300-at valamiről, ami árt a pénztárcájuknak.

	      Emellett más szabályokat is be lehet vezetni, hogy a rendszer gyorsan működjön. Például szabhatnánk
	      egy 7 napos határidőt minden javaslatra. Ha az emberek nem jutnak konszenzusra ez idő alatt, a képviselő
	      beletása szerint szavazhat.

	      Azonban az emberek a kezükbe akarják venni az irányítást, és erre nincs jobb annál, mint amikor mi
	      magunk alkotjuk a szabályokat, amelyek szerint élünk. Ha a demokráciában való részvételt könnyűvé
	      tesszük, az emberek részt fognak venni. Az internet erre kiváló eszköz. Végül is ha valamire, arra
	      kiváló, hogy az emberek kinyilvánítsák a véleményüket rajta.

	\item \textbf{Civilek nem tudnak törvényt alkotni.} A politikusok sem. A törvényeket többnyire
	      jogászok, köztisztviselők, tanácsadók, cégek, lobbicsoportok vagy más érintettek írják a politikus
	      körül. Sok országban további lépések biztosítják, hogy a megfogalmazás világos legyen és megfeleljen
	      az  alkotmánynak. A lényeg, hogy a törvényeket ritkán írja egyetlen ember --- és ez így van jól.
	      Civilek is együtt tudnak dolgozni a segíteni hajlandó szakértőkkel, és így a jogszabályt addig
	      lehet formálni, amíg jól megfogalmazott és elfogadható nem lesz.

	      Valamilyen módon hatnunk kell a kormányra és minden tisztelettel, a tüntetések nem működnek. Olyan
	      emberekre van szükség, akik leülnek, lépésről lépésre kidolgozzák a jogszabályokat, gondoskodnak róla,
	      hogy a változás hatékony és igazságos legyen, majd keresztülviszik azt a politikusokon és
	      a milliárdosokon. Ez a rendszer ezt lehetővé teszi ezt.
\end{enumerate}

Ezzel az új rendszerrel olyan lehetőségek tárulnak fel az átlag polgár előtt, ami eddig a kivételesek
számára volt csak elérhető. Humanista csoportok tucatjai kezdhetnek saját előterjesztéseket írni és ha
a többség kegyét megnyerik törvényeket alkotni belőlük. A következő lista olyan előterjesztéseket
tartalmaz, amik érdekesek lehetnek az átlag polgár számára és jó eséllyel a többségi támogatást is
megszereznék. Azonban ez a mozgalom nem támogatja ezeket hivatalosan, azon egyszerű oknál fogva, hogy
a direkt demokrácia egyetlen célja, hogy erőt adjon az átlag ember kezébe és nem az, hogy megmondja
mi a teendő ezzel az erővel. Ennek ellenére, ezek az ötletek érdekesek lehetnek, ha másért
nem is azért, hogy elindítson egy beszélgetést arról milyen változásokra lenne szükség az országban:

\begin{itemize}

	\item \textbf{Otthoni Munkavégzés}\cite{work_from_home}: a Covid világjárvány megmutatta, hogy a
	      legtöbb irodai munka elvégezhető távolról. Ez a törvény nem kerül semmibe a kormánynak mégis
	      rengeteg, azonnal érezhető előnye van:
	      \begin{enumerate}
		      \item Javítja a városok levegőminőségét, mivel kevesebb embernek kell ingáznia.

		      \item Enyhíti a lakhatási válságot, mert az emberek elköltözhetnek a városból anélkül,
		            hogy elveszítenék az állásukat.

		      \item Csökkenti a zajt a nagyvárosokban, mivel kevesebb járműre van szükség a mindennapokban.

		      \item Támogatja az autómentes infrastruktúrát, hiszen kevesebb autó kevesebb helyet foglal.

		      \item A parkolóhelyek iránti kisebb igény felszabadít területeket lakóházak, parkok vagy más
		            városi szükségletek számára.

		      \item Még azok is profitálnak belőle, akik továbbra is bejárnak: alacsonyabb bérleti díjak,
		            kevesebb forgalom és tisztább levegő várja őket.

		      \item Életrekelti az elnéptelenedő kisvárosokat, mivel azon dolgozók akik csendesebb életre
		            vágynak elköltözhetnek a fővárosból anélkül, hogy elvesztenék a munkájukat.
	      \end{enumerate}

	\item \textbf{4 Napos Munkahét}: ahogy korábban említettük egy 32 órás munkahétnek is rengeteg előnye
	      lenne és ennek megvalósításához sem szükséges számottevő állami erőforrás:
	      \begin{enumerate}
		      \item Csökkenti a munkahelyi stresszt.

		      \item Egyes esetekben javítja a munkahelyi tejlesítményt.

		      \item Több időt biztosít a polgároknak saját érdeklődéseikre.
	      \end{enumerate}

	\item \textbf{Lakhatási és Közlekedési Hivatal}: különösen sokat tudna segíteni a fiataloknak
	      vagy a lakhatással küzdőknek. Ez egy lakhatási és közlekedési ügynökség létrehozása lenne. Ez
	      a kormányzati ügynökség ott építene lakásokat és tömegközlekedést, ahol szükség van rájuk, egészen
	      addig, amíg mindenki, aki szeretne, tud bérelni. A bérleti díjat egyetlen bérlő fizetésének 25\%-ában
	      kellene meghatározni --- jóval a 35\%-os EU-us ajánlás alatt, és messze a majdnem 50\% alatt, amit
	      sokan ma fizetnek. A kormány a beszedett bérleti díjat felhasználhatná a közösségi lakások fenntartására
	      és bővítésére. Ezeket a lakóprojekteket tömegközlekedéssel kell összekötni, ezért az ügynökségnek
	      felhatalmazást kellene kapnia közlekedési infrastruktúra építésére is. Létfontosságú, hogy az ügynökség
	      ezeket a projekteket közvetlenül irányítsa, ne pedig alvállalkozók seregén keresztül, mivel utóbbi esetben
	      az ár többszörösét kell majd az ügynökségnek kifizetnie. Bécs bebizonyította mennyi előnnyel jár, ha a
	      kormány segít a polgárainak megélni, de itt van néhány:
	      \begin{enumerate}
		      \item Csökkenti a piacon lévő albérletek árát és a lakbéreket.

		      \item Az átlag embernek több jövedelme marad, amit kedve szerint költhet el, ami segíti a gazdaságot.

		      \item A kormánynak lehetősége lesz ténylegesen segíteni az embereknek egy gazdasági válság vagy a
		            következő világjárvány esetén. Például úgy, hogy nem gyűjti be lakbért azoktól, akik elvesztették
		            a munkájukat a válság alatt.

		      \item Csökkenti az egyenlőtlenségeket a társadalomban, azáltal, hogy a kormány talpra tudja állítani
		            a rászorulókat. Például biztosíthat lakhatást hajléktalanoknak, amíg munkát találnak, vagy otthonról
		            elzavart LMBTQ fiataloknak.
	      \end{enumerate}
\end{itemize}

\subsection{Útmutatás Valódi Rendszerváltáshoz}

A fakultatív modern közvetlen demokrácia szépsége, hogy ellentétben sok egyéb ötlettel ami jól hangzik elméletben,
de nincs egyértelmű út a megvalósításához, ezt viszonylag könnyű elérni. Ugyan rengeteg polgárt meg kell
győzni, elég őket arról meggyőzni, hogy vegyék saját kezükben az irányítást ahelyett, hogy átruháznák
azt a következő négy évre. Ennél fogva, a választók egy széles köre megszólítható. A terv néhány pontban
leírható:
\begin{enumerate}
	\item A lehető legtöbb ember meggyőzése --- a politikai nézetektől függetlenül --- egy adott
	      választókerületben, majd kapcsolatfelvétel egy megbízható képviselővel vagy jelöltel. Ha nincs
	      ilyen egy független delegálása. A választott egyén egyetlen dolga, hogy a készre kapott
	      törvénymódosításra szavazzon, ami lehetővé teszi majd a közvetlen demokráciát.

	\item Ismételjük meg ezt minden választókerületre.

\end{enumerate}

Azonban, hogy a változás ténylegesen megtörténhessen szükség van önkéntesekre. Jogászokra, média szakértőkre,
vagy egy olyan csoporthoz tartozókra, akiket eddig elhnyagolt a politika. Részletek ezen a linken:
https://www.facebook.com/groups/1429998928614117/permalink/1430000145280662/